\subsection{Weyl involutions reconstruct the $36$ double-sixes, the $45$ tritangents, and their cross-incidence}
\label{sec:20260917-e6-involution-36-45}

There is a second, independent route to the cubic-surface $27$--$36$--$45$
triangle that does not start from the W33 apartment/code carrier.  Work in the
Picard lattice of a smooth cubic surface with basis $H,E_1,\ldots,E_6$ and
intersection form $\operatorname{diag}(1,-1^6)$.  The $27$ $(-1)$-curve
classes are
\[
 E_i,\qquad H-E_i-E_j,\qquad
 2H-\sum_{j\ne i}E_j.
\]
The six simple root reflections generate the full Weyl permutation group
\[
 |W(E_6)|=51840,
\]
and the sign kernel has order $25920$.

Exhaustive enumeration of the $891$ nonidentity involutions gives the four
class sizes
\[
 \boxed{36,\ 270,\ 540,\ 45},
\]
corresponding to involution degrees $1,2,3,4$.  The simultaneous fixed-carrier
fingerprints on the $(27,36,45)$ objects are
\[
\begin{array}{c|c|c|c|c}
 d & |C_d| & \#\operatorname{Fix}(27) &
 \#\operatorname{Fix}(45) & \#\operatorname{Fix}(36)\\ \hline
 1&36 &15&15&16\\
 2&270& 7& 5& 8\\
 3&540& 3& 7& 4\\
 4&45 & 3&13&12.
\end{array}
\]
The degree census is classical; the objectwise identifications below are the
new executable reconstruction.

\paragraph{The $36$ reflection class is literally the double-six carrier.}
For every root reflection $r$, the $12$ cubic lines moved by $r$ are exactly the
support of one double-six.  The $36$ moved supports are distinct and exhaust the
$36$ double-sixes.  Thus
\[
 \boxed{\{\text{reflections of }W(E_6)\}\ \longleftrightarrow\
        \{\text{double-sixes}\}}.
\]
Equivalently, the classical root-pair/double-six correspondence is visible
directly in the $27$-line permutation action, without importing labels.

\paragraph{The $540$ odd degree-three involutions form a $12$-sheeted cover of the $45$ tritangents.}
Every degree-three involution fixes exactly three of the $27$ lines pointwise,
and those three lines form a tritangent.  Conversely every tritangent occurs as
the fixed triple of exactly twelve such involutions.  Hence
\[
 \boxed{540=45\cdot12}.
\]
Write $F_T$ for the twelve degree-three involutions above a tritangent $T$,
and let $r_D$ be the reflection corresponding to a double-six $D$.

\paragraph{The group law reconstructs the $45\times36$ incidence matrix.}
A direct census gives the sharp dichotomy
\[
 |T\cap D|=0
 \quad\Longleftrightarrow\quad
 \#\{w\in F_T:wr_D=r_Dw\}=4,
\]
while
\[
 |T\cap D|=2
 \quad\Longleftrightarrow\quad
 \#\{w\in F_T:wr_D=r_Dw\}=0.
\]
Therefore the commutation relation internal to $W(E_6)$ alone reconstructs the
tritangent/double-six cross-incidence.  The exact census is
\[
 \boxed{0^{540}+2^{1080}},
 \qquad 45\cdot36=540+1080.
\]
This is precisely the cross-incidence independently reconstructed in
Pass~4659 from the W33 internal $27/36$ carrier, where
$T^{\mathsf T}R=2(J-D)$ and the same $0^{540},2^{1080}$ intersection census
appears.  Thus the W33/code route and the classical Weyl-involution route meet
objectwise, not merely numerically.

The executable certificate is
\texttt{analysis/w33\_e6\_involution\_reconstructs\_double\_six\_tritangent.py},
with frozen output
\texttt{data/w33\_e6\_involution\_reconstructs\_double\_six\_tritangent.json}.
The classical involution-degree census agrees with Carter's Weyl-group
classification and with the cubic-surface involution discussion of
Dolgachev--Duncan.  The finite theorem does not assert a physical $E_6$ gauge
field; transport to the W33/Suzuki carriers uses the separately certified
$27$--$36$--$45$ equivariant intertwiners.
